\section{怎样分析多重复句}

如果一个复句由三个以上的分句构成，而且有两个以上的层次，那么，这样的复句就叫做多重复句。如果分句虽是三个以上，但只有一个层次的，那也不算多重复句。例如：

（1）在中华民族的开化史上，有素称发达的农业和

（并列）

手工业①，$\vert$ 有许多伟大的思想家、科学家、发明家、

\hspace{7em}（并列）

政治家、军事家和文艺家②，
$\vert$ 有丰富的文化典籍③。

\hspace{17em}（因果）

（2）因为我们是为人民服务的①，$\vert$所以，我们如果

\hspace{3em}（假设）

有缺点②，$\Vert$就不怕别人批评指出③。

这两个复句都由三个分句构成。但是，第一个复句的\ 
①\ 
②\ 
③\ 
三个分句之间的关系都是并列关系，也就是说，整个复
句只有一层并列关系，所以不是多重复句。而第二个复句，
就全句而言，是一个因果句，①是因，②③是果，这是第一
层关系；②③是假设关系，这是第二层关系：因此全句是个
多重复句。

我们怎样才能正确地分析多重复句呢？

\textbf{1. 必须要熟悉复句的各种关系。}

对于什么样的复句，是什么关系，经常用哪些关联词语
来表示，我们必须要熟悉。有些复句没有关联词语，这就要
从分句的内容上看它们之间的联系，有时还要借助上下文，
来判断分句之间的关系。根据全句所表达的意思和常用的关
联词语，能判断出分句间是一种什么关系，这是分析多重复
句的基础和条件。

一般说来，经常成对地使用的关联词语有如下这些：

表示并列关系：也······也······；又······又······；一方面
······一方面······；有时······有时······；一会儿······一会儿······
······；既······又······，等等。（单独使用的有：也、又、同
时、同样、同时也、同时又······）

表示连贯关系：起先······后来······；首先······然后······
（单独使用的有：就、于是、便、才、后来、然后、接着、
跟着、又······。有时不用关联词语，只是依照时间或空间的
顺序而连续。）

表示递进关系：不但······而且······；不仅······还······；不但不（不但没有）······反而······，连······也······；尚且······何况······；别说······连······也······，等等。（单独使用的有：而且、并且、甚至······）

表示选择关系：不是······就是······；要么······要么······；要就是······要就是······；是······还是······；与其······不如······；宁可······也不······；宁愿······不愿······；或者······或者······，等等，（单独使用的有：还是、或、或者、或是、不然、否则······）

表示转折关系：虽然······但是······；尽管······却······；虽（虽是）······但（但是，可是，而）······等等。（单独使用的有：但是、但、可是、然而、而、却、倒、反、偏······）

表示假设关系：如果······就······；倘若（倘使、若）······便（就，那，那么，那么就，那就）······；要是（要······的话，要不是）······那么······；即使······也······；纵然（就算，就是，哪怕）······也（还）······，等等。（单独使用的有：就、便、那么、那么就、那就······）

表示条件关系：只要······就······；只有······才······；除非······才（不）······；无论······都······；不论（不管、任、任凭）······总（总是、也）······，等等，（单独使用的有：就、才、都、总、总是、也······）

表示因果关系：因为······所以······；由于······因而······；既然······那么······，等等。（单独使用的有：所以、因此、因而、那么、就······）

表示目的关系的关联词语一般都单独使用，有“以、以便、好、为的是、为了、省得、免得、以免”等。这种复句的一个分句表示行为，而有表示目的、关系的、关联词语的分句则表示行为的目的。

\textbf{2. 要准确地划清多重复句中的分句。}

把复句中的各个分句准确地划分出来，确定这个多重复句究竟包含几个分句，是分析多重复句的第一步工作。为此，就该注意：

（1）不要把作为分句的句子成分的主谓词组、动宾词组、联合词组等误作为分句，因为分句在全句结构中都是相对独立的。

（2）几个分句共有一个主语的，后面的分句主语一般承前省略；由于表达的需要，分句主语有时也有蒙后省略的。我们要把省略主语的分句看成是复句中的一个分句，而不能误把它当成一个句子成分看待。

（3）分句与分句之间有语音停顿，书面上常用逗号或分号隔开，有时用冒号隔开。

（4）一般说来，多重复句的分句之间，常有互相呼应的关联词语。只是要留心，有些关联词语是连接词与词、词组与词组或句子与句子的，不要把这些关联词语误当作连接分句与分句的。

\textbf{3. 总观全句，首先找出第一个层次的关系。}

第一层次是整个多重复句的主要层次。分析多重复句，首先要根据关联词语，结合全句的意思，划出第一个层次的前后几个部分，判定第一层次的关系。如果是没有关联词语的“意合句”（由意义紧密相关而组合起来的复句），能根据前后分句之间的意义补上关联词语的，可以补上去判定它们的关系。第一层次的关系抓准了，全句的基本结构也就搞清楚了，对于以下层次的分析就有了前提条件。这一点，跟对单句内部划分主、谓两大部分的要求有些类似。大层次划错了，小层次就无法划分，或者勉强划分也不对头。第一层次划得对不对，是可以检验出来的。例如：

（1）不管反动派怎样企图阻止历史车轮的前进①，
\noindent
（条件）\hspace{9em}（递进）

$\vert$革命或迟或早总会发生②，$\Vert$并且将必然取得胜利③。

\hspace{15em}（转折）

（2）我们不需要死记硬背①，$\vert$但是我们需要用基本

\hspace{14em}（因果）

\noindent
事实的知识来发展和增进每个学习者的思考力②，$\Vert$因为

\hspace{17em}（假设）

\noindent
不把学到的全部知识融会贯通③，$\Vert$$\vert$共产主义就会变成空中

\hspace{1em}（并列）\hspace{8em}（并列）

\noindent
楼阁④，$\Vert$$\Vert$$\vert$就会成为一块空招牌⑤，$\Vert$$\Vert$共产主义者也只会是一些吹牛家⑥。

我们先来看第一个多重复句，它包括①②③三个分句，用了以下的关联词语：“不管······总······”、“并且”。“不\\
管······总······”表示条件关系，“并且”表示递进关系。总观全句，分句①提出了一个条件，分句②③表示在①句所提出的条件下会产生的结果。这样划分，不但各部分都能说得上一定的关系，而且各部分加起来就包括句子的全体，没有丢下任何一个部分，这就说明划得对。如果划得不对，下一层就划分不下去。因为如果把第一个层次划在“并且”之前的话，那么，全句则成了递进关系的联合复句，而与③句构成递进关系的只是②句，①句则被丢掉了，这当然就划错了。

我们再来看第二个多重复句，它包括①②③④⑤⑥六个分句，用了“但是”、“因为”、“就”、“也只”等关联词语。那么第一层该从哪儿划分呢?有人可能认为应该在①②跟③--⑥之间划开，前一部分是结果，后一部分是原因。可是仔细一看，因果关系只表现在②跟③--⑥之间。这种划法则把①丢掉了。可见这种划法是不对的。正确的划法应该是：第一层在①跟②--⑥之间，全句是个表示转折关系的偏正复句。这样划分，不仅各部分都能说得上一定的关系，而且各部分加起来就包括了句子的全体，没有丢下任何一个部分，这就说明划对了。

当然，一个多重复句如果在分句之间用分号隔开，用分号的前后两个部分一定是第一层关系，而且一般是并列关系，这是没有疑问的。

\textbf{从大到小，逐层分析。}
在弄清第一层次的基础上，进一步分析第一层次前后各个部分的内部结构，逐层理清第二、第三······层次的结构关系，直到不可分为止，这样全句的结构就能逐步分析清楚。

例如，上面第三点中的第一个多重复句，我们总观全句，确定了第一个层次在①和②两个分句之间，它们的意义关系是条件关系。在这之后，我们再分别分析第一层次前后各个部分的内部结构及其意义关系。①句之前已没有别的分句，即是说①句已是最小的分句，我们就不再分析；②句之后还有③句，那么我们就根据“并且”这个表示递进关系的关联词语，可以立即确定第二个层次就划在②和③两个分句之间，它们是表示递进关系的联合复句。分析到此，全句没有更小的分句了，说明已经分析完毕。

再如，上面第三点中的第二个多重复句，我们总观全句，确定了第一个层次在①和②两个分句之间，它们的意义关系是转折关系。①句已分析到最小分句，不再分析；②句之后，还有四个分句，我们再由大到小、逐层分析下去。那就是：第二层在②跟③——⑥之间，是因果关系；第三层在③跟④⑤⑥之间，是假设关系；第四层在⑤跟⑥之间，是并列关系；第五层在④跟⑤之间，是并列关系。

总之，只要第一层关系的前后两部分，仍然包含两个或两个以上的分句，就得一层一层地分析下去，找出分句和分句之间的意义关系，用“$\vert$”、“$\Vert$”、“ $\Vert$$\vert$”······逐层划开，标明各层关系的意义，直至不可再分为止。我们可以再看三个例句，从中好好体会这种由大到小、逐层分析的方法：

\hspace{10em}（因果）

（1）我赞美白杨树，$\vert$ 就因为它不但象征了北方的农

\noindent
（递进）

\noindent
民，$\Vert$ 尤其象征了今天我们民族解放斗争中所不可缺的朴质、坚强、力求上进的精神。

\hspace{20em}（递进）

（2）我们不但要有大批杰出的科学技术专家，$\Vert$ 而且

\hspace{9em}（目的）

\noindent
必须动员亿万人民刻苦学习，$\Vert$$\vert$ 提高全民族科学文化水平，

\noindent
（条件）\hspace{12em}（并列）

$\vert$ 才能在各个领域赶超世界先进水平，$\Vert$ 才能加快四个现代化的进程。

（3）我们要和一切资本主义国家的无产阶级联合起来，
\newpage
（并列）

$\Vert$要和日本的、英国的、美国的、德国的、意大利的以及

（条件）

一切资本主义国家的无产阶级联合起来，$\vert$才能打倒帝国主义

\hspace{2em}（并列）\hspace{8em}（并列）

义，$\Vert$ 解放我们的民族和人民，$\Vert$$\vert$解放世界的民族和人民。

第一句是个二重复句，共有三个分句；第三句是个三重复句，共有五个分句；第三句也是三重复句，共有五个分句。我们准确地确定了全句的第一层关系之后，然后再将第一层关系的前后两个部分一层一层地分析下去，直到分句都是单句不可再分为止。不论一个多重复句层次如何多，关系如何复杂，只要掌握了分析多重复句的方法，再作一些有关的练习，就一定能够准确无误、得心应手地将它分析清楚。
\newpage
